Reciprocal-space construction of the TAPW basis for the representative
zero-twist hBN/R2G supercell. (a) First Brillouin zones of graphene (blue) and hBN
(red). The gray tiling shows repeated moir\'e Brillouin zones, and the filled gray
hexagon at $\Gamma$ marks the central moir\'e Brillouin zone. The box marks the
selected valley, enlarged in (b,c). (b) The 27 retained hBN reciprocal vectors
$\widetilde{\mathbf G}$ (red circles). (c) The 37 reciprocal vectors retained for
each graphene layer (blue circles). Diamonds mark the material-specific valley
centers $\mathbf K_{\rm hBN}$ and $\mathbf K_{\rm Gr}$. Both enlarged panels use
the same origin $\mathbf K_{\rm hBN}$ and momentum scale $|\mathbf b_{M,1}|$.
Dashed circles indicate radial shells about the corresponding material's valley
center, with $\mathbf Q_l=\mathbf K_l-\widetilde{\mathbf G}$. The first six hBN
shells contain 3, 3, 6, 6, 3, and 6 vectors, while the first six graphene shells
contain 1, 6, 6, 6, 12, and 6 vectors; the first graphene shell is the central
$\mathbf Q=0$ point. Each vector carries two sublattice orbital labels (B/N for
hBN and A/B for graphene). With one hBN layer, two graphene layers, and spin
neglected, the TAPW basis contains $2\times27+2\times2\times37=202$ states per valley.
